%% kapa: exam=1 round=37 year=2026 period=2 form=A subject=accounting subject_ko=회계학 seq=5 number=45
%% source: Q-Net 감정평가사 기출문제 (cst00309) artlSeq=5258924 file=raw/1008/5258924/2026년 제37회 감정평가사 1차 2교시.pdf
%% conversion: kordoc 4.24.51 -> tools/gen1.py (auto); review=proofread-sonnet; flags=none
\begin{question}{45}
20×1년 초 설립한 (주)감평의 20×1년 기말 상품 관련 자료는 다음과 같다. (주)감평이 항목별기준 저가법과, 조별기준 저가법을 각각 적용할 경우 20×1년 포괄손익계산서에 인식할 재고자산평가손실은?
\par
{\small\centering
\begin{tabular}{|c|c|c|c|c|c|}\hline
조 & 항목 & \makecell{장부\\수량} & \makecell{실제\\수량} & \makecell{단위당\\원가} & \makecell{단위당\\순실현가능가치} \\\hline
\multirow{2}{*}{A} & A1 & 100개 & 80개 & ￦1,000 & ￦1,100 \\\cline{2-6}
 & A2 & 100개 & 100개 & ￦1,200 & ￦1,150 \\\hline
\multirow{2}{*}{B} & B1 & 200개 & 150개 & ￦2,000 & ￦1,900 \\\cline{2-6}
 & B2 & 200개 & 200개 & ￦2,100 & ￦2,050 \\\hline
\end{tabular}\par}
\hspace*{2.4em}\makebox[4.5cm][l]{\uline{항목별기준}}\uline{조별기준} \par
\begin{choices}
  \item \makebox[4.5cm][l]{￦30,000,}￦22,000
  \item \makebox[4.5cm][l]{￦30,000,}￦25,000
  \item \makebox[4.5cm][l]{￦30,000,}￦30,000
  \item \makebox[4.5cm][l]{￦35,000,}￦25,000
  \item \makebox[4.5cm][l]{￦35,000,}￦30,000
\end{choices}
\end{question}
\answer{2}
